\chapter{为什么症状不能直接说明原因}

\begin{statebox}
\textbf{项目 v1.0 的证据基线。} 当前 73 项回归把单元、集成、固定种子随机、
ELF/镜像审计、QEMU 正常路径和 22 条历史失败路径组合起来。109 个键盘字符、
四进程统计、管道守恒、文件系统双启动和损坏拒绝分别证明不同边界；任何一项都不能
由一条 \code{READY} 日志替代。正文继续说明真实生产系统怎样把这些局部证据组织成
可定位延迟与故障的观测体系。
\end{statebox}

\inputchapterbody{chapters/ch17-panic-observability}

\section{调试工具与性能边界}
\inputdetail{chapters/ch22-observability-debugging}

\begin{problemframe}
\textbf{全书得到的系统。} 从无操作系统机器开始建立的每一项能力，如今都有明确对象、
所有者、入口、等待关系、提交点、失败路径和观察证据。书末综合案例重新追踪一次文件保存，
但此时其中每个对象都已在前文分别推导，因此完整链路只承担综合验证，不再承担术语开场。
\end{problemframe}
